\chapter{Results and Discussion}

\section{Data Characteristics}

The feature-selection analysis was evaluated across two registered Gene Expression Omnibus (GEO) datasets after log2 transformation and $z$-score normalization:
\begin{itemize}
    \item \textbf{GSE19804 (Lung Cancer):} 120 samples (60 Cancer, 60 Normal), 54,675 probes. Balanced class distribution.
    \item \textbf{GSE42568 (Breast Cancer):} 121 samples (104 Cancer, 17 Normal), 54,675 probes. Severe class imbalance (~6:1 ratio).
\end{itemize}
Both datasets exhibit extreme high dimensionality with a probe-to-sample ratio exceeding 450:1.

\section{Standardized Feature-Selection Outcomes}

Under the updated pipeline implementation, feature selection was executed across both datasets using a standardized target subset size of $k = 20$ features per method. Filter methods applied a $p$-value constraint ($p \leq 0.05$) combined with top-$k$ significance ranking, while wrapper methods utilized an ANOVA prefilter (top 500) before SVM-RFE or Random Forest ranking.

Table~\ref{tab:feature_selection_counts} summarizes the selection outcomes across all five implemented algorithms.

\begin{table}[H]
\centering
\caption{Feature selection dimensionality reduction summary ($p \le 0.05, k = 20$).}
\label{tab:feature_selection_counts}
\begin{tabular}{llcc}
\toprule
\textbf{Category} & \textbf{Method Key} & \textbf{Selected Probes} & \textbf{Reduction (\%)} \\
\midrule
Filter & \texttt{filter\_ttest} (Welch's $t$-test) & 20 & 99.96\% \\
Filter & \texttt{filter\_anova} (ANOVA $F$-test) & 20 & 99.96\% \\
Wrapper & \texttt{wrapper\_svm} (SVM-RFE) & 20 & 99.96\% \\
Wrapper & \texttt{wrapper\_rf} (RF Importance) & 20 & 99.96\% \\
Embedded & \texttt{embedded\_lasso} (LASSO $L_1$) & 20 & 99.96\% \\
\bottomrule
\end{tabular}
\end{table}

\section{Feature Selection Patterns on GSE19804}

On the balanced GSE19804 lung cancer dataset, all 54,675 gene probes were evaluated. Statistical significance filtering ($p \le 0.05$) identified 19,505 significant probes for Welch's $t$-test and 19,540 probes for the ANOVA $F$-test. Ranking these significant probes yielded the top 20 most discriminative probes for each filter method.

Top probes identified by univariate filters included highly significant lung cancer biomarkers such as \texttt{204446_s_at} ($p < 10^{-15}$) and \texttt{209189_at}. The wrapper methods (SVM-RFE and Random Forest) identified complementary multivariate probe sets, while LASSO selected a sparse linear combination of probes.

\section{Feature Selection Patterns on GSE42568}

On the imbalanced GSE42568 breast cancer dataset, 54,589 non-constant probes were evaluated. Despite the severe class imbalance (104 cancer vs 17 normal), univariate filter methods and wrapper methods successfully identified 20 statistically significant probes ($p < 0.05$). The selected feature sets effectively reduced the feature space by 99.96\%, isolating highly informative transcriptional probes suitable for downstream classifier training.

\section{Comparative Discussion}

Standardizing the target feature size to $k = 20$ across Filter, Wrapper, and Embedded methods provides a consistent basis for comparing algorithm behavior:
\begin{enumerate}
    \item \textbf{Filter Methods:} Computationally efficient ($< 1$ second runtime), selecting univariate statistical leaders.
    \item \textbf{Wrapper Methods (SVM-RFE \& RF):} Capture feature-feature interaction dependencies, providing optimized subsets tailored for linear margin classifiers.
    \item \textbf{Embedded Methods (LASSO):} Provide intrinsic regularization, forcing weak coefficients to zero while maintaining sparsity.
\end{enumerate}

All selected feature lists, p-value statistics, and gene symbol annotations have been saved in the project repository under \texttt{results/feature_selection/}.